\documentclass[11pt,a4paper]{article}

\usepackage[margin=2.5cm]{geometry}
\usepackage{amsmath, amssymb}
\usepackage{graphicx}
\usepackage{physics}
\usepackage{enumitem}

\title{Project 1 -- Airplanes}
\date{February 25, 2026}

\begin{document}

\maketitle

\noindent
\textbf{Deadline:} 11.3.2026, 13:00.

\vspace{0.5em}

Please hand in your solution as follows: One single pdf containing all text and plots, and one single, working (!)
python script or notebook containing all code used. Ideally, give your documents an easily identifiable name such as
``Project1\_groupA\_text.pdf''.

Your submission does not need to win a beauty contest, but make sure that it is clear:
readable plots with clear figure captions and references in the text, clear section titles, clear language etc.
A very messy hand-in may loose you up to 1 out of 10.

Recall that if you use ChatGPT or similar tools, you must briefly explain how you did it and how you checked that its
results make sense, and hand in a complete log sheet with your prompts and the answers obtained.
This will not explicitly be graded, but serves as a reference.

To be handed in via BB, one per group. All names of group members should be on the sheet.

\section*{Part 1 -- A flat airplane wing (analytic solution)}

In the analytic part we look at a flat airplane wing. As we will see, it generates lift by being turned by an angle $\alpha$
(the angle of attack) with respect to the wind, which is horizontal. For simplicity, however, we rotate our coordinate
system such that the wing is along the $x$-axis, meaning that the wind comes slightly from ``below'' in our coordinate
system, even though physically speaking, it is the wind that is horizontal.

Consider a flat plate of length $4b$, situated horizontally (i.e.\ from $x=-2b$ to $x=2b$).
We know that we can consider this plate as product of a Joukowski-transformation of a cylinder with radius $1$
(centred on the origin).

The circle is parameterised by its angle:
\[
Z(\gamma) = R e^{i\gamma}
\]
where here $R=1$ (radius). In the following, $Z$ is the complex coordinate in the situation with a cylindrical wing
(circle in 2D), and $z$ its Joukowski transform.

The Joukowski transformation is given by
\[
z = Z + \frac{b^2}{Z}
\]
and we will take $b=1$.

Now let’s say that the plate is in a wind field, whose complex velocity far away from the origin takes the value
\[
U_x - i U_y = U e^{-i\alpha}
\]
where $0 \le \alpha \le \frac{\pi}{2}$ is the angle of attack, and $U$ is real positive.

\subsection*{a) (5 points)}

Explain briefly why the complex potential around the cylinder is given by
\[
f(Z) = U e^{-i\alpha} Z + \frac{U e^{i\alpha}}{Z}
- \frac{i\Gamma}{2\pi} \ln(Z)
\tag{1}
\]
with an as yet unknown parameter $\Gamma$.

If you use equations from the lecture notes, refer to them explicitly (e.g., ``starting from eq.~1.1, we see...'').

\subsection*{b) (BONUS: max +5 points)}

We have seen in the lecture that the complex potential in point $z$ near the plate equals that of point $Z$ near the cylinder,
where $z$ is the Joukowski transform of $Z$.
This could give us the idea that to obtain the velocity potential at any $z$ near the plate, we just need to seek the
corresponding value $Z$ and check the velocity potential there.

Compute the inverse of the Joukowski transformation and show that it is not unique.
Argue that for our case of a flat horizontal plate (and $b=1$), it makes sense to pick always the solution $Z$ that has
the same sign of the imaginary part as $z$, i.e.\ $Y$ and $y$ have the same sign.
Come up with a rule for points on the $x$-axis with $|z|>2$.

\subsection*{c) (12 points)}

We are interested in the velocity distribution on the plate, in particular the edges ($z=-2$, $z=2$).
Show that the complex velocity around the plate,
\[
\frac{df}{dz},
\]
fulfils
\[
\frac{df}{dz}
=
\frac{U \sin(\gamma - \alpha) - \Gamma/(4\pi)}
{\sin(\gamma)}.
\]

Hint: Use
\[
\frac{df}{dz} = \frac{df/dZ}{dz/dZ}
\]
and first show that this relation indeed holds.

\subsection*{d) (8 points)}

What is the problem at the tips of the plate?

Hint: First check which angles $\gamma$ correspond to the tips of the plate.

\subsection*{e) (15 points)}

For the trailing edge (the back edge, turned away from the wind), we can fix the problem by a smart choice of $\Gamma$.
For the leading edge, we cannot do the same simultaneously.

In proper wings (as opposed to flat plates), this problem is solved by making the wing slightly rounded at the leading
edge, while the trailing edge remains sharp. So we consider it sufficient here to remove the singularity on the trailing edge.

\begin{itemize}[leftmargin=1.5em]
\item What is the correct choice of $\Gamma$?
\item Compute the drag and lift force using the Blasius theorem and the theorem of Kutta–Joukowski.
\item Discuss its dependence on the angle of attack.
\item Compute the unit of the lift force in our equation. Why is it not Newton? What does this imply for the lift of an actual plane?
\end{itemize}

\section*{Part 2 -- A realistic wing profile (numerical)}

In the second part, we investigate the lift on a more realistically-shaped airplane wing using again the Joukowski transformation.
We will put $b=1$ throughout.

We start with a cylinder in the $Z$-plane:
\[
Z = Z_0 + R e^{i\gamma}, \qquad 0 \le \gamma < 2\pi.
\]

(Notation: $z$ is the coordinate in the airfoil case; $Z$ in the cylinder case.)

\subsection*{a) (6 points)}

Plot the Joukowski transform of a circle with
\[
Z_0 = -0.1 + 0.22i, \qquad R = 1.12.
\]

You should get a profile with a sharp edge on the right and rounded on the left.

\subsection*{b) (8 points)}

Plot the streamfunction around the cylinder for
\[
\Gamma = 0 \quad \text{and} \quad \Gamma = -3.
\]

Check the boundary condition carefully.

Plot the velocity potential as well. Why is there a weird jump along some (probably horizontal) line?
Does it have a physical meaning?

\subsection*{c) (8 points)}

Using the Joukowski transformation, also plot the stream function around the wing
for the same two values of $\Gamma$.
Check the boundary condition.

\subsection*{d) (6 points)}

What happens to the stagnation point around the sharp trailing edge?
If needed, zoom in.

\subsection*{e) (6 points)}

Find a value of the circulation for which the flow around the trailing edge is smooth.
(You may determine it numerically to percent precision.)
Use this value from now on.

\subsection*{f) (10 points)}

For this circulation, compute and plot the pressure field around the cylinder and the wing.

Hints:
\[
w = \frac{df}{dz} \approx \frac{f_1 - f_2}{z_1 - z_2}
\]

Think carefully about possible artefacts in the velocity field.

\subsection*{g) (8 points)}

Numerically compute the lift using
\[
F_x - iF_y
=
\frac{i\rho}{2}
\oint_C w^2 \, dz
\approx
\frac{i\rho}{2}
\sum_k w_k^2 \Delta z_k,
\qquad
\Delta z_k = \frac{z_{k+1} - z_{k-1}}{2}.
\]

\subsection*{h) (8 points)}

Does the result in g) match the theoretical results from the lectures?

\vspace{1em}

\noindent
\textbf{Technical hints:}
\begin{itemize}
\item Useful plotting tools can be found in the python introduction code.
\item \texttt{plt.contour} can generate contour plots even on irregular grids.
\end{itemize}

\end{document}